\section{Related Work}
\label{sec:related-work}

\paragraph{KV memory systems.}
PagedAttention and SGLang provide page and prefix management inside an LLM
runtime~\cite{kwon2023vllm,zheng2024sglang}.  Mooncake and HiCache extend KV
storage across host memory and distributed backends
~\cite{qin2025mooncake,xie2025hicache}; SYMPHONY, Strata, and Bidaw further
optimize disaggregation and hierarchical placement
~\cite{agarwal2026symphony,xie2026strata,hu2026bidaw}.  \system{} uses these
mechanisms as storage and transfer substrates.  Its focus is the workflow value
used to choose which failure-recoverable state occupies their limited fast
capacity.

\paragraph{Agent-aware serving.}
Parrot, KVFlow, Pythia, and Agentix expose application structure or
predictability to scheduling and cache reuse
~\cite{lin2024parrot,pan2025kvflow,yu2026pythia,luo2026agentix}.  TokenCake,
Continuum, PBKV, and TokenDance use criticality, waiting time, predicted reuse,
or synchronized execution to manage KV for agent workloads
~\cite{bian2026tokencake,li2026continuum,zheng2026pbkv,bian2026tokendance}.
These systems establish that workflow information benefits normal serving.
\system{} instead asks how one dynamic Slack value should govern both
pre-failure protection and post-failure recovery.

\paragraph{Fault-tolerant LLM serving.}
D\'ej\`aVu streams KV state for recovery, while GhostServe and RaidServe protect
or reconstruct serving state with GPU/host redundancy
~\cite{strati2024dejavu,jayakody2026ghostserve,xu2026raidserve}.  KevlarFlow,
LUMEN, Concordia, AnchorTP, Tarragon, and ReviveMoE address complementary
recovery, checkpointing, or reconfiguration problems
~\cite{qian2026kevlarflow,cao2026lumen,gan2026concordia,xu2026anchortp,
zhang2026tarragon,li2026revivemoe}.  CacheFlow accelerates the physical
restore/load/recompute path~\cite{nian2026cacheflow}.  \system{} does not claim
a new copy engine or general reconfiguration mechanism; it supplies the online
workflow priority and KV-management policy that select and order such work.

\paragraph{Classical mechanisms.}
Landlord provides a credit/rent view of online variable-size caching
~\cite{young2002online}, and real-time systems have long used Slack to place
checkpoints and schedule recovery~\cite{melhem2004interplay,acharya2008dynamic}.
\system{} does not claim either idea in isolation.  \textsc{SlackLandlord}
specializes the credit to continuous KV recovery harm and reuses the same Slack
after a serving failure; \textsc{GuardedAIMD} adds a slow feedback envelope
around optional peer-HBM capacity while local allocation remains authoritative.
